% -----------------------------------------------------------------------------
% Conteudo AODV-EN. Template base: Prof. Wesley Pacheco Calixto (IFG).
% -----------------------------------------------------------------------------
\chapter{INTRODUÇÃO}\label{introduuxe7uxe3o}

No início da década de 1990, Mark Weiser, pesquisador do Xerox Palo Alto
Research Center (PARC), introduziu em seu artigo seminal \emph{The
Computer for the 21st Century} o conceito de Computação Ubíqua, ao
afirmar que as tecnologias mais profundas são aquelas que desaparecem,
integrando-se de forma invisível ao cotidiano humano \cite{weiser1991}.
Essa visão pioneira antecipou um ecossistema composto por centenas de
dispositivos computacionais de baixo custo, baixo consumo energético e
amplamente interconectados por redes sem fio, operando em diferentes
escalas espaciais e comunicando-se de forma transparente entre si e com
o ambiente. Mais de três décadas depois, tal paradigma materializa-se
por meio da Internet das Coisas (\emph{Internet of Things} -- IoT) e das
Redes de Sensores Sem Fio (\emph{Wireless Sensor Networks} -- WSNs),
tecnologias que conectam bilhões de dispositivos embarcados mundialmente
em aplicações que vão desde automação residencial e agricultura de
precisão até monitoramento ambiental e sistemas industriais inteligentes
\cite{priyadarshi2025}.

Nesse contexto, o microcontrolador ESP32, desenvolvido pela Espressif
Systems, destaca-se como uma plataforma versátil e amplamente adotada,
combinando conectividade Wi-Fi e Bluetooth integradas, elevado poder de
processamento e custo acessível, características que o tornam ideal para
prototipagem e implantação de soluções IoT em escala. Entretanto, a
comunicação em redes IoT frequentemente enfrenta desafios relacionados à
cobertura, escalabilidade e resiliência. Topologias centralizadas, como
a estrela, apresentam limitações significativas: a dependência de um
ponto de acesso único constitui um gargalo de desempenho e um ponto
único de falha, comprometendo a robustez do sistema \cite{chaizeng2021}.

Como alternativa, as Redes Mesh Sem Fio (\emph{Wireless Mesh Networks}
-- WMNs) emergem como uma solução promissora, oferecendo características
de auto-organização (\emph{self-configuration}), auto-recuperação
(\emph{self-healing}) e comunicação ponto-a-ponto entre os nós
participantes. Nas redes mesh, a comunicação em múltiplos saltos
(\emph{multi-hop}) permite que pacotes de dados sejam encaminhados
através de nós intermediários até alcançarem seu destino final,
estendendo significativamente o alcance da rede sem a necessidade de
aumentar a potência de transmissão dos dispositivos individuais
\cite{arregui2025}. Essa abordagem possibilita a implantação de redes em
áreas extensas utilizando dispositivos de baixa potência, reduzindo
custos e consumo energético. Em contrapartida, a comunicação multi-hop
introduz desafios adicionais relacionados ao roteamento de pacotes,
gerenciamento de interferência e manutenção da topologia da rede,
exigindo algoritmos especializados para garantir a entrega confiável dos
dados.

Para enfrentar esses desafios, os protocolos de roteamento para redes
ad-hoc podem ser classificados em três categorias principais: proativos,
reativos e híbridos. Os protocolos proativos, como o \emph{Optimized
Link State Routing} (OLSR), mantêm tabelas de roteamento constantemente
atualizadas, oferecendo baixa latência na descoberta de rotas às custas
de maior \emph{overhead} de controle. Os protocolos reativos, como o
\emph{Ad-hoc On-Demand Distance Vector} (AODV), descobrem rotas apenas
quando necessário, reduzindo o tráfego de controle em redes com tráfego
esporádico. Já os protocolos híbridos combinam características de ambas
as abordagens, adaptando-se às diferentes condições da rede
\cite{chaizeng2021}.

Entre os mecanismos de comunicação disponíveis no ecossistema ESP32, o
protocolo ESP-NOW destaca-se como uma alternativa de alto desempenho
para comunicação direta entre dispositivos. O ESP-NOW opera de forma
\emph{connectionless} sobre a camada
de enlace do padrão IEEE 802.11 e permite a troca de mensagens curtas
(até 250 bytes de \emph{payload}) sem a necessidade de estabelecimento
prévio de conexão ou associação a um ponto de acesso Wi-Fi
\cite{becker2025}. Estudos experimentais recentes demonstram as
vantagens do ESP-NOW em relação a outras tecnologias de comunicação sem
fio. \citeonline{becker2025}, em experimentos de campo conduzidos em
ambientes abertos e florestais, reportaram latências médias de 2,8 ms e
\emph{Packet Delivery Ratio} (PDR) superior a 99\% em distâncias de até
55 metros em condições de linha de visada. \citeonline{urazayev2023},
avaliando o protocolo em ambiente interno, constataram que o ESP-NOW
apresenta alcance 15\% superior e consumo energético 30\% inferior
quando comparado a soluções baseadas em Wi-Fi TCP convencional. Essas
características tornam o ESP-NOW especialmente atrativo para aplicações
IoT em ambientes remotos ou alimentados por baterias de capacidade
limitada.

Apesar de suas vantagens, o protocolo apresenta limitações estruturais
que restringem sua aplicabilidade em cenários que demandam redes de
maior escala ou alcance estendido. A primeira limitação refere-se à
ausência de suporte nativo a roteamento multi-hop: o ESP-NOW foi
concebido exclusivamente para comunicações ponto-a-ponto ou em topologia
estrela, não oferecendo mecanismos intrínsecos para encaminhamento de
pacotes através de nós intermediários \cite{becker2025}. A segunda
limitação diz respeito ao número máximo de \emph{peers} simultâneos
suportados: cada dispositivo ESP32 pode manter no máximo 20 \emph{peers}
registrados em sua tabela de pares, o que impõe restrições severas à
escalabilidade da rede. A terceira limitação relaciona-se à ausência de
mecanismo de \emph{broadcast} nativo eficiente, dificultando a
disseminação de mensagens de controle necessárias para protocolos de
descoberta de rotas. Essa lacuna torna-se evidente ao se analisarem
soluções de redes em malha amplamente utilizadas no ecossistema
Espressif, como o ESP-WIFI-MESH e a biblioteca PainlessMesh, que operam
sobre a pilha Wi-Fi convencional (IEEE 802.11). Conforme discutido por
\citeonline{sestak2022}, o uso do Wi-Fi padrão para manutenção da
topologia mesh impõe elevado \emph{overhead} de processamento e consumo
energético, inviabilizando sua aplicação em dispositivos estritamente
alimentados por baterias e comprometendo a principal vantagem
competitiva do ESP-NOW. Trabalhos recentes têm buscado preencher essa
lacuna: \citeonline{cujilema2023} propuseram o algoritmo BRAM-NOW para
criação de redes mesh seguras sobre ESP-NOW, demonstrando latência média
de 75 ms e taxa de perda de pacotes de 9,25\% em ambiente residencial.
Entretanto, tal abordagem não se fundamenta em protocolos de roteamento
padronizados, o que dificulta sua comparação com soluções estabelecidas
na literatura e limita sua extensibilidade.

Diante desse cenário, formula-se a seguinte questão de pesquisa: é
viável adaptar um protocolo de roteamento padronizado e amplamente
validado, como o AODV, para operar sobre o ESP-NOW, contornando suas
restrições de \emph{peers}, ausência de \emph{broadcast} confiável e de
suporte a multi-hop, mantendo desempenho competitivo frente a uma
abordagem de referência baseada em inundação (\emph{flooding})? Como
hipótese central, admite-se que um roteamento reativo sob demanda,
adaptado a essas restrições, é capaz de entregar pacotes em múltiplos
saltos sobre ESP-NOW com taxa de entrega comparável à da inundação
controlada, porém a um custo de controle e de ocupação de canal
significativamente menor à medida que a rede cresce em tamanho e
diâmetro.

Para responder a essa questão, o presente trabalho propõe o
desenvolvimento e a avaliação de um algoritmo de roteamento denominado
AODV-EN (\emph{Ad-hoc On-Demand Distance Vector for ESP-NOW}), uma
adaptação do protocolo AODV, especificado na RFC 3561 \cite{perkins2003}.
A escolha do AODV como protocolo base justifica-se por suas
características de roteamento reativo sob demanda, que minimizam o
tráfego de controle em redes com comunicação intermitente, e por sua
ampla documentação e validação na literatura acadêmica, facilitando
comparações e extensões futuras. As principais adaptações propostas para
adequar o AODV às restrições do ESP-NOW incluem: (i) gerenciamento
dinâmico de \emph{peers} utilizando política LRU (\emph{Least Recently
Used}), com cache configurável de entradas ativas, permitindo que o
protocolo alcance nós em número superior ao do cache ao longo do tempo,
dentro do limite físico de 20 \emph{peers} do ESP-NOW; (ii) disseminação
das mensagens RREQ (\emph{Route Request}) por \emph{flooding} controlado
sobre \emph{broadcast} de enlace do ESP-NOW (FF:FF:FF:FF:FF:FF), com
supressão de duplicatas por cache de identificadores e limite de
propagação (TTL), preservando a semântica clássica do AODV; e (iii)
incorporação de métrica híbrida de roteamento que combina contagem de
saltos (\emph{hop count}) com indicador de qualidade do sinal recebido
(RSSI), possibilitando a seleção de rotas que equilibrem menor distância
lógica e melhor qualidade de enlace.

\section{OBJETIVOS}\label{objetivos}

Esta seção apresenta o objetivo geral e os objetivos específicos que
orientam o desenvolvimento deste trabalho.

\subsection{Objetivo geral}\label{objetivo-geral}

Propor, implementar e avaliar um algoritmo de roteamento multi-hop
baseado no protocolo AODV, adaptado para operação em redes mesh
construídas sobre o protocolo ESP-NOW, considerando as restrições
específicas dessa tecnologia e utilizando hardware de baixo custo.

\subsection{Objetivos específicos}\label{objetivos-especuxedficos}

Para alcançar o objetivo geral, foram definidos os seguintes objetivos
específicos:

\begin{enumerate}
\def\labelenumi{\alph{enumi})}
\tightlist
\item
  Analisar as limitações técnicas do protocolo ESP-NOW para comunicação
  multi-hop, identificando os desafios específicos relacionados ao
  limite de \emph{peers}, ausência de \emph{broadcast} nativo e
  gerenciamento de rotas;
\item
  Projetar as adaptações necessárias ao protocolo AODV para adequá-lo às
  restrições do ESP-NOW, incluindo gerenciamento de \emph{peers} com
  política LRU, \emph{flooding} controlado e métrica híbrida de
  roteamento;
\item
  Implementar um protótipo funcional do algoritmo AODV-EN em módulos
  ESP32, desenvolvendo as estruturas de dados, mensagens de controle e
  mecanismos de descoberta e manutenção de rotas;
\item
  Avaliar o desempenho do algoritmo proposto através de experimentos em
  diferentes topologias de rede, utilizando métricas quantitativas como
  \emph{Packet Delivery Ratio} (PDR), latência fim-a-fim,
  \emph{Normalized Routing Load} (NRL) e consumo energético;
\item
  Comparar os resultados obtidos com trabalhos correlatos da literatura,
  posicionando a contribuição do AODV-EN em relação às soluções
  existentes para redes mesh sobre ESP-NOW e outras tecnologias.
\end{enumerate}

\section{CENÁRIOS DE APLICAÇÃO}\label{cenarios-de-aplicacao}

A motivação para prover roteamento multi-hop padronizado sobre ESP-NOW
torna-se concreta em cenários nos quais a infraestrutura de rede é
ausente, intermitente ou economicamente inviável, mas a comunicação entre
muitos dispositivos de baixo custo e baixo consumo é necessária ao longo
de áreas que excedem o alcance de um único enlace. Quatro classes de
aplicação ilustram a relevância da proposta:

\begin{itemize}
\item \textbf{Agricultura de precisão e monitoramento rural.} Sensores de
  umidade do solo, temperatura e nível de reservatórios distribuídos por
  uma propriedade extensa, sem cobertura de Wi-Fi ou celular, podem
  encaminhar leituras até um concentrador por múltiplos saltos, a custo de
  hardware da ordem de poucos dólares por nó.
\item \textbf{Monitoramento industrial e predial.} Em galpões, plantas
  fabris e edifícios, paredes e estruturas metálicas fragmentam a
  conectividade; uma malha auto-organizável que contorna obstáculos por
  rotas alternativas é mais robusta que enlaces ponto a ponto.
\item \textbf{Redes comunitárias e resposta a desastres.} Em regiões sem
  infraestrutura ou após eventos que a destroem, nós alimentados por
  bateria podem formar rapidamente uma malha temporária para troca de
  mensagens e telemetria, sem depender de pontos de acesso.
\item \textbf{Automação residencial de baixo custo.} Dispositivos
  domésticos (sensores, atuadores) distribuídos por uma casa de vários
  cômodos, onde a atenuação por paredes já exige multi-hop, conforme
  observado no próprio cenário experimental deste trabalho.
\end{itemize}

Em todos esses contextos, a adoção de um protocolo \emph{padronizado}
(AODV, RFC 3561), em vez de uma solução \emph{ad hoc} proprietária,
favorece a interoperabilidade, a reprodutibilidade e a manutenção de longo
prazo --- diferencial central da proposta.

\section{ESTRUTURA DO TRABALHO}\label{estrutura-do-trabalho}

Além desta introdução, este trabalho está organizado em mais cinco
capítulos. O Capítulo~\ref{cap:referencial} apresenta o referencial
teórico, abordando as redes de sensores sem fio e redes mesh, o
protocolo ESP-NOW e suas limitações, os protocolos de roteamento ad-hoc
e o protocolo AODV, o algoritmo de inundação utilizado como referência e
as métricas de avaliação de desempenho. O Capítulo~\ref{cap:metodologia}
descreve a metodologia adotada, fundamentada na \emph{Design Science
Research}, detalhando os cenários experimentais, os parâmetros de
configuração e o procedimento de coleta e análise dos dados. O
Capítulo~\ref{cap:implementacao} apresenta o projeto e a implementação
do AODV-EN, incluindo a arquitetura do firmware, as estruturas de dados,
o formato das mensagens e as adaptações propostas. O
Capítulo~\ref{cap:resultados} reporta e discute os resultados obtidos na
avaliação experimental em \emph{hardware} e em simulação, comparando o
AODV-EN com o algoritmo de referência e com a literatura correlata. Por
fim, o Capítulo~\ref{cap:conclusao} apresenta as conclusões, as
limitações do estudo e as propostas de trabalhos futuros.
